% \akari{need to make sure to cite concurrent work on idea generations, automating data science / machine learning coding ... etc }
\paragraph{Scientific LMs.}
% Previous research has explored the potential of LLMs to accelerate scientific discovery. 
% For example, \textsc{Galactica}~\citep{taylor2022galactica}, a scientific language model trained on vast text data from diverse sources such as academic papers and knowledge bases, aimed to encapsulate extensive scientific knowledge to aid discovery. 
Scientific LMs have spanned various domains, including biomedical \citep{SciFive2021,BioBART2022,BioGPT2022}, medical \citep{singhal2023large,MeLlama2024,PMCLlama2023,Meditron2023}, biomedical \citep{BioMedGPT2023,BioMedGPT2024,BioMistral2024}, geoscience \citep{K22023}, astronomy~\citep{AstroLLaMA2023} and multidisciplinary science \citep{DARWIN2023}, with some models such as SciGLM~\citep{SciGLM2024} and UniSmart~\citep{UniSmart2024} that aim to cover diverse scientific domains in a single model. 
Recently, several works show that powerful general-purpose LLMs can also show strong capabilities in scientific tasks, such as medical question answering~\citep{ai4science2023impact,singhal2023large}, chemistry experimentation~\citep{zheng2023gpt} and applied mechanics~\citep{10.1115/1.4062773}. 
However, the language model's reliance on information memorized within its parameters leads to frequent hallucinations in its output~\citep{li-etal-2024-dawn}. 

\paragraph{LMs to assist scientists.}
Recent studies have also examined LLMs' capabilities to assist scientists in performing a range of scientific procedures, including generating novel research ideas~\citep{baek2024researchagent,Yang2023LargeLM} and automating experimental code generation~\citep{Huang2023MLAgentBenchEL,tian2024scicode}. Our work, however, focuses specifically on benchmarking and developing methods for automating literature reviews and addressing questions related to up-to-date research—tasks that are crucial to, and particularly challenging, for scientific inquiry. 
Several concurrent studies have attempted to build retrieval-augmented pipelines using proprietary LLMs and external APIs (e.g., Semantic Scholar API) for scientific literature review agents~\citep{agarwal2024litllm,skarlinski2024language,wang2024autosurvey}. While these studies and our research all explore the potential of retrieval-augmented LMs in automating literature synthesis, prior works often relied on proprietary, black-box systems and limited evaluations, which commonly entail small-scale human evaluation or simplified setups such as multiple-choice QA. 
In contrast, our work introduces a comprehensive benchmark with automated metrics, involves user studies with experts across three scientific disciplines, and develops new methodologies to train specialized open models. \model significantly outperforms previously introduced systems and shows superiority over human experts in five domains. 

% \paragraph{LLMs that accelerate science.}
% \gncomment{Also, do we only want to focus on LLMs for science? There's also pre-LLM work on various aspects, including survey generation \citep{jha-etal-2015-content}, etc. so maybe that should be fleshed out.}
\paragraph{Benchmarks for scientific literature understanding.}
Several works have developed benchmarks to evaluate models' abilities to understand scientific literature. Prior datasets, such as SciFact~\citep{wadden-etal-2020-fact}, QASPER~\citep{dasigi-etal-2021-dataset}, and QASA~\citep{lee2023qasa}, largely focus on single-paper settings, where the necessary information to answer queries is contained within a single pre-selected paper. However, in real-world scenarios, experts often need to synthesize information from multiple papers to answer questions. To address this gap, \data introduces newly annotated tasks that require reasoning across multiple papers.
There are also scientific summarization tasks, such as Multi-XScience~\citep{lu2020multi}, where models are provided with multiple papers and asked to generate summaries, typically based on the related work sections of those papers. However, in this work, we focus on scenarios where the relevant papers are not specified in advance, making the task more challenging. 
Recently, \citet{xu-etal-2024-kiwi} introduced KIWI, a dataset containing 200 questions and human-verified or edited answers generated by state-of-the-art LLMs, with a focus on the NLP domain. KIWI also provides a set of relevant papers that models must consider. While both KIWI and \data feature multi-paper, information-seeking tasks, \data includes both human-written answers and automatic evaluation pipelines. In contrast, KIWI focuses more on human evaluations, and its reference answers are primarily model-generated.

% \paragraph{General-domain long-form evaluations. }



% Moreover, we conducted human evaluations across four scientific domains, comparing system-generated responses to expert-written answers for open-ended literature review questions.

